\section{Epoch-wise Double Descent}
\label{sec:epoch-dd}

%In this section, we provide evidence for the generality of our notion of EMC by studying the interaction between training epochs and test error. We present a novel form of double-descent with respect to training epochs.


%In this section, we study the interaction between train epochs, model size, and test error.

%First, we consider the evolution of test error of a fixed, large architecture as training progresses.
%Increasing train epochs increases effective model complexity in our framework -- and thus a sufficiently large model transitions from underparameterized to overparameterized over the course of training.
% Indeed, consistent with our generalized double-descent hypothesis, we observe
% an ``epoch-wise'' double descent for large models:
% test error first decreases (at the beginning of training), then increases (around the interpolation threshold), then decreases once more (at the end of training).


In this section, we demonstrate a novel form of double-descent with respect to training epochs, which is consistent with our unified view of effective model complexity (EMC) and the generalized double descent hypothesis.
Increasing the train time increases the EMC---and thus a sufficiently large model transitions from under- to over-parameterized over the course of training.

\begin{figure}[h]
    \centering
    \begin{subfigure}[t]{.48\textwidth}
        \centering
        \includegraphics[width=\textwidth]{epochDD_regimes}
    \end{subfigure}
    \begin{subfigure}[t]{.48\textwidth}
        \centering
        \includegraphics[width=\textwidth]{epochDD_ocean_plot}
    \end{subfigure}
    \caption{
    {\bf Left: } 
    Training dynamics for models in three regimes.
    Models are ResNet18s on CIFAR10 with 20\% label noise,
    trained using Adam with learning rate $0.0001$, and data augmentation.
    {\bf Right:} Test error over (Model size $\x$ Epochs).
    Three slices of this plot are shown on the left.
    }
    \label{fig:epochDD-three-regimes}
\end{figure}

As illustrated in Figure~\ref{fig:epochDD-three-regimes}, 
sufficiently large models can
undergo
a ``double descent'' behavior where test error first decreases then increases near the interpolation
threshold, and then decreases again.
In contrast, for ``medium sized'' models, for which training to completion will only barely reach $\approx 0$ error, the
test error as a function of training time will follow a classical U-like curve where it is better to stop early.
Models that are too small to reach the approximation threshold will remain in the ``under parameterized'' regime where increasing
train time monotonically decreases test error.
Our experiments
(Figure~\ref{fig:epochDD-noise-dataset-arch}) show that many settings
of dataset and architecture exhibit epoch-wise double descent, in the presence of label noise.
Further, this phenomenon
is robust across optimizer variations and learning rate schedules (see additional experiments in Appendix \ref{subsec:appendix_epoch}).
As in model-wise double descent, the test error peak is accentuated with label noise. 

\iffalse

Indeed, we observe three distinct train-time behaviors
depending on the model size and label noise (Figure~\ref{fig:epochDD-three-regimes}):
\begin{enumerate}
    \item Small models (which are under-parameterized throughout training):
    Test error decreases monotonically throughout training. 
    \item Intermediate-size models
    (which transition from under- to critically parameterized over training): Test error decreases but starts diverging close to interpolation (conventional understanding of when over-fitting sets in).
    \item Large models (which transition from under- to over-parameterized throughout training):
    Test error decreases, increases or plateaus near the interpolation threshold, then decreases again.\todo{added "or plateaus"--Boaz}
\end{enumerate}

\fi

%These behaviors are consistent with our generalized double descent hypothesis.

%We demonstrate this phenomenon for a variety of natural models, datasets, and optimizers.
%Epoch-wise double-descent is most prominent in settings where model-wise double-descent
%is prominent.

% This manifests most prominently as a ``double-descent'' in test error in settings with label noise -- here,

% Here, we provide further evidence of a double-descent in test error with the number of epochs. 

\begin{figure}[h]
    \centering
    \begin{subfigure}[t]{.32\textwidth}
        \centering
        \includegraphics[width=0.99\linewidth]{epochDD_cf10_labelnoise_train_test}
        \caption{ResNet18 on CIFAR10.}
    \end{subfigure}
    \begin{subfigure}[t]{.32\textwidth}
        \centering
        \includegraphics[width=0.99\linewidth]{epochDD_cf100_labelnoise_train_test}
        \caption{ResNet18 on CIFAR100.}
    \end{subfigure}
    \begin{subfigure}[t]{.32\textwidth}
        \centering
        \includegraphics[width=0.99\linewidth]{epochDD_cf10_mcnn_labelnoise_train_test}
        \caption{5-layer CNN on CIFAR 10.}
    \end{subfigure}
    \caption{{\bf Epoch-wise double descent} for ResNet18 and CNN (width=128).
    ResNets trained using Adam with learning rate $0.0001$,
    and CNNs trained with SGD with inverse-squareroot learning rate.
    }
    \label{fig:epochDD-noise-dataset-arch}
\end{figure}


Conventional wisdom suggests that training is split into two phases: (1) In the first phase, the network learns a function with a small generalization gap (2) In the second phase, the network starts to over-fit the data leading to an increase in test error.
Our experiments suggest that this is not the complete picture---in some regimes, the test error decreases again and may achieve a lower value at the end of training as compared to the first minimum (see Fig \ref{fig:epochDD-noise-dataset-arch} for 10\% label noise).

% A practical suggestion from this is that you don't want to stop just when test error starts to diverge.


% Conventional wisdom is that when you train, you start to over-fit at some point (where you should early stop). 
